(1) 均匀弦小振幅横波方程为 \(\mu y_{tt}=T y_{xx}\)，波速 \(c=\sqrt{T/\mu}\)。固定端位移为零，\(y(0,t)=0\)；无横向端部载荷的自由端满足 \(Ty_x(L,t)=0\)，所以 \(y_x(L,t)=0\)，这与第二个固定端完全不同。
(2) 设 \(y(x,t)=Y(x)e^{i\omega t}\)，空间方程 \(Y''+k^2Y=0\)，\(k=\omega/c\)。左端固定使 \(Y(x)=A\sin(kx)\)，右端自由使 \(Y'(L)=Ak\cos(kL)=0\)。非零解要求 \(k_nL=(n+1/2)\pi\)，\(n=0,1,\ldots\)。
(3) 由 \(k_n=(n+1/2)\pi/L\)，得到 \(\omega_n=(n+1/2)\pi c/L\)。基模角频率 \(\omega_0=\pi c/(2L)\)，频率 \(f_0=c/(4L)\)，波长 \(\lambda_0=c/f_0=4L\)；两端固定弦基频为 \(c/(2L)\)，恰是该混合边界基频的两倍。
(4) 机械波能量 \(E=\frac12\int_0^L[\mu y_t^2+T y_x^2]dx\)。乘波方程以 \(y_t\) 后积分分部，得 \(\dot E=[Ty_xy_t]_0^L\)。左端因固定而 \(y_t(0)=0\)，右端因自由而 \(y_x(L)=0\)，两项都为零，故 \(\dot E=0\)。
